\documentclass[11pt]{article}

% Change "review" to "final" to generate the final (sometimes called camera-ready) version.
% Change to "preprint" to generate a non-anonymous version with page numbers.
\usepackage[review]{acl}

% Standard package includes
\usepackage{times}
\usepackage{latexsym}

% For proper rendering and hyphenation of words containing Latin characters (including in bib files)
\usepackage[T1]{fontenc}
% For Vietnamese characters
% \usepackage[T5]{fontenc}
% See https://www.latex-project.org/help/documentation/encguide.pdf for other character sets

% This assumes your files are encoded as UTF8
\usepackage[utf8]{inputenc}

% This is not strictly necessary, and may be commented out,
% but it will improve the layout of the manuscript,
% and will typically save some space.
\usepackage{microtype}

% This is also not strictly necessary, and may be commented out.
% However, it will improve the aesthetics of text in
% the typewriter font.
\usepackage{inconsolata}

%Including images in your LaTeX document requires adding
%additional package(s)
\usepackage{graphicx}

% If the title and author information does not fit in the area allocated, uncomment the following
%
%\setlength\titlebox{<dim>}
%
% and set <dim> to something 5cm or larger.

\title{Severity-Aware Conformal Risk Control for Medical LLM Hallucinations}

% Author information can be set in various styles:
% For several authors from the same institution:
% \author{Author 1 \and ... \and Author n \\
%         Address line \\ ... \\ Address line}
% if the names do not fit well on one line use
%         Author 1 \\ {\bf Author 2} \\ ... \\ {\bf Author n} \\
% For authors from different institutions:
% \author{Author 1 \\ Address line \\  ... \\ Address line
%         \And  ... \And
%         Author n \\ Address line \\ ... \\ Address line}
% To start a separate ``row'' of authors use \AND, as in
% \author{Author 1 \\ Address line \\  ... \\ Address line
%         \AND
%         Author 2 \\ Address line \\ ... \\ Address line \And
%         Author 3 \\ Address line \\ ... \\ Address line}

\author{First Author \\
  Affiliation / Address line 1 \\
  Affiliation / Address line 2 \\
  Affiliation / Address line 3 \\
  \texttt{email@domain} \\\And
  Second Author \\
  Affiliation / Address line 1 \\
  Affiliation / Address line 2 \\
  Affiliation / Address line 3 \\
  \texttt{email@domain} \\}

%\author{
%  \textbf{First Author\textsuperscript{1}},
%  \textbf{Second Author\textsuperscript{1,2}},
%  \textbf{Third T. Author\textsuperscript{1}},
%  \textbf{Fourth Author\textsuperscript{1}},
%\\
%  \textbf{Fifth Author\textsuperscript{1,2}},
%  \textbf{Sixth Author\textsuperscript{1}},
%  \textbf{Seventh Author\textsuperscript{1}},
%  \textbf{Eighth Author \textsuperscript{1,2,3,4}},
%\\
%  \textbf{Ninth Author\textsuperscript{1}},
%  \textbf{Tenth Author\textsuperscript{1}},
%  \textbf{Eleventh E. Author\textsuperscript{1,2,3,4,5}},
%  \textbf{Twelfth Author\textsuperscript{1}},
%\\
%  \textbf{Thirteenth Author\textsuperscript{3}},
%  \textbf{Fourteenth F. Author\textsuperscript{2,4}},
%  \textbf{Fifteenth Author\textsuperscript{1}},
%  \textbf{Sixteenth Author\textsuperscript{1}},
%\\
%  \textbf{Seventeenth S. Author\textsuperscript{4,5}},
%  \textbf{Eighteenth Author\textsuperscript{3,4}},
%  \textbf{Nineteenth N. Author\textsuperscript{2,5}},
%  \textbf{Twentieth Author\textsuperscript{1}}
%\\
%\\
%  \textsuperscript{1}Affiliation 1,
%  \textsuperscript{2}Affiliation 2,
%  \textsuperscript{3}Affiliation 3,
%  \textsuperscript{4}Affiliation 4,
%  \textsuperscript{5}Affiliation 5
%\\
%  \small{
%    \textbf{Correspondence:} \href{mailto:email@domain}{email@domain}
%  }
%}

\begin{document}
\maketitle
\begin{abstract}
This document is a supplement to the general instructions for *ACL authors. It contains instructions for using the \LaTeX{} style files for ACL conferences.
The document itself conforms to its own specifications, and is therefore an example of what your manuscript should look like.
These instructions should be used both for papers submitted for review and for final versions of accepted papers.
\end{abstract}

\section{Introduction}
Large language models (LLMs) are increasingly being used to support patients and healthcare professionals by answering medical questions and assisting with health-related information. Yet they state false claims as confidently as true ones. A model might report that a drug should be taken at "5000 mg per day" just as fluently as it reports that "metformin treats type 2 diabetes," with no signal that one of them is wrong. In most applications, factual errors are often minor and carry limited consequences. In medicine, a wrong dose or a missed contraindication can cause real harm. For these systems to be trusted in the clinic, it is not enough to be usually right: we need a reliable bound on how often, and how dangerously, they are wrong.

Raw confidence estimates from LLMs do not provide such a guarantee. Both softmax probabilities and self-reported confidence scores are often poorly calibrated, making them unreliable indicators of the true probability of error. Conformal prediction offers a way to obtain this guarantee regardless of how well-calibrated the underlying confidence score is. Given a held-out calibration set of examples with known correctness, it converts any confidence score into a decision threshold with a distribution-free guarantee. On new data, the error rate among kept predictions stays below a target level α that the user chooses in advance.



Existing conformal factuality guarantees overlook a basic asymmetry: not all hallucinations are equally harmful. A single budget, however, treats every hallucination as an identical unit
of error. Because the guarantee constrains only the combined rate, a filter can be highly accurate on benign claims and still have a very high error rate on dangerous claims, as long as the benign claims, which are far more numerous, pull the average back within budget. This is not hypothetical. Across two open models and three medical QA benchmarks, a single global CRC threshold tuned to a 10% marginal budget leaves the dangerous-claim hallucination rate between 0.068 and 0.094, well above a 0.05 dangerous target, on every model and dataset pair we test. 

We address this by making the guarantee severity-aware. Rather than controlling one harm-blind error rate, we partition claims into the two severity tiers above and assign each tier its own conformal budget. Because the calibration procedure enforces the budgets separately, holding the dangerous tier to a stricter target than the benign tier directly forces a higher confidence threshold on dangerous claims, rather than relying on severity to come out in the wash of an aggregate rate. This reallocates strictness rather than raising it everywhere: the looser benign budget lets the system surface more true benign information, even as the stricter dangerous
budget catches more dangerous hallucinations. 


These instructions are for authors submitting papers to *ACL conferences using \LaTeX. They are not self-contained. All authors must follow the general instructions for *ACL proceedings,\footnote{\url{http://acl-org.github.io/ACLPUB/formatting.html}} and this document contains additional instructions for the \LaTeX{} style files.

The templates include the \LaTeX{} source of this document (\texttt{acl\_latex.tex}),
the \LaTeX{} style file used to format it (\texttt{acl.sty}),
an ACL bibliography style (\texttt{acl\_natbib.bst}),
an example bibliography (\texttt{custom.bib}),
and the bibliography for the ACL Anthology (\texttt{anthology.bib}).

\section{Related Work}
\paragraph{Conformal prediction and risk control for LLMs.}
Conformal prediction calibrates a decision threshold on labeled data to obtain
distribution-free, finite-sample guarantees
\citep{vovk2005algorithmic,angelopoulos2023gentle}. Conformal risk control
(CRC) generalizes the guarantee from coverage to any bounded loss that
decreases as the threshold tightens \citep{angelopoulos2024crc}. Applied to
language models, this underlies conformal language modeling
\citep{quach2024conformal}, conformal factuality, which removes low-confidence
sub-claims to bound the retained-error rate \citep{mohri2024conformal}, and
conformal abstention, which declines to answer when uncertainty is high
\citep{abbasiyadkori2024conformal}. A recent line examines how the realized
risk of CRC varies across calibration splits and uses randomization and
bootstrapping to characterize it beyond its mean \citep{pang2025bootstrap}; we
adopt a related cluster bootstrap to report confidence intervals and per-split
violation rates. All of these methods control a single, harm-blind error rate.

\paragraph{Claim-level and medical factuality.}
Verifying long-form answers one atomic claim at a time follows the
decompose-then-verify approach of FactScore \citep{min2023factscore} and SAFE
\citep{wei2024safe}; MedScore adapts the decomposition step to clinical text,
where general-purpose decomposers underperform \citep{huang2025medscore}. Most
relevant to us, \citet{cherian2024llm} bring claim-level conformal factuality to
medical long-form QA with conditional, topic-adaptive guarantees. We build on
this claim-level medical conformal factuality; what we change is the quantity
being controlled, not the decomposition or scoring pipeline. For claim-level
labels we use K-QA \citep{manes2024kqa} and subsets of MedLFQA.

\paragraph{Harm-aware factuality and safety.}
A growing body of work observes that not all claims matter equally.
Importance-sensitive metrics weight factual errors by their relevance to the
query \citep{wanner2025importance}, but they are evaluation measures and carry
no statistical guarantee. On the guarantee side, CARE \citep{bedi2026care} is a
conformal safety layer for medical summarization that controls separate budgets
for hallucination and omission at the sentence level; its losses are unweighted,
and it explicitly leaves harm-weighted objectives and type-specific risk budgets
to future work, citing the lack of validated severity weights. Our work provides
exactly that missing piece: a separate conformal budget per clinical severity
tier, with the tiers validated against physician labels, at the claim level for
medical question answering.

Severity-conditional control also relates to Mondrian and group-conditional
conformal prediction \citep{vovk2012conditional} and to weighted CRC, which
reweights within a single budget \citep{xu2024conformal}. To our knowledge, we
are the first to assign each clinical severity tier its own conformal budget for
claim-level medical factuality, turning the harm-blind guarantee of prior
conformal factuality into a harm-stratified one.


\section{Method}
\label{sec:method}

\subsection{Problem setup and notation}
\label{sec:setup}
We work at the level of atomic claims. A model answer to a medical question is
decomposed into claims $c_1, \dots, c_m$ following the decompose-then-verify
recipe \citep{min2023factscore, mohri2024conformal}. Each claim $c$ carries three
quantities: a confidence score $s(c) \in [0,1]$, an interpretation of the model's
estimated probability that the claim is true; a binary label $y(c) \in \{0, 1\}$,
where $y=1$ marks a hallucination and $y=0$ a faithful claim; and a severity tier
$t(c) \in \mathcal{T} = \{\textsc{benign}, \textsc{dangerous}\}$ assigned by a
clinical grader. The label is used only during calibration; at deployment a claim
is filtered using its score and tier alone.

The filter is a single threshold rule. Given a threshold $\lambda \in [0,1]$, we
\emph{keep} claim $c$ when $s(c) \ge \lambda$ and \emph{flag} it otherwise.
Because we cannot in general verify a flagged claim, the object we control is the
rate at which kept claims are hallucinations. For a claim $c$ we define the
per-claim loss
\begin{equation}
\ell(c; \lambda) \;=\; \mathbf{1}\!\left[s(c) \ge \lambda\right]\,
\mathbf{1}\!\left[y(c) = 1\right],
\label{eq:loss}
\end{equation}
the indicator that a hallucination survives the filter. This loss is bounded in
$[0,1]$ and is monotone non-increasing in $\lambda$: raising the threshold only
removes claims, so it can only turn a kept hallucination into a flagged one, never
the reverse. These two properties are exactly what conformal risk control requires.

\subsection{Conformal risk control}
\label{sec:crc}
Let $\{c_i\}_{i=1}^{n}$ be a calibration set of claims with known labels. The
empirical risk of the threshold $\lambda$ is the mean loss
\begin{equation}
\widehat{R}_n(\lambda) \;=\; \frac{1}{n} \sum_{i=1}^{n} \ell(c_i; \lambda)
\;=\; \frac{1}{n} \sum_{i=1}^{n}
\mathbf{1}\!\left[s(c_i) \ge \lambda\right]\,\mathbf{1}\!\left[y(c_i) = 1\right],
\label{eq:emp-risk}
\end{equation}
that is, the fraction of \emph{all} calibration claims that are both kept and
hallucinated. Conformal risk control \citep{angelopoulos2024crc} selects the
smallest threshold whose finite-sample-corrected risk meets a target budget
$\alpha$,
\begin{equation}
\widehat{\lambda} \;=\; \inf\left\{ \lambda :\;
\frac{n}{n+1}\,\widehat{R}_n(\lambda) + \frac{B}{n+1} \;\le\; \alpha \right\},
\label{eq:crc-threshold}
\end{equation}
where $B = 1$ upper-bounds the loss in \eqref{eq:loss}. Choosing the smallest such
$\lambda$ is deliberate: it is the least aggressive filter that still meets the
budget, so it keeps as many claims as possible. Under exchangeability of the
calibration claims and a fresh test claim $c_{n+1}$, this threshold satisfies the
distribution-free guarantee
\begin{equation}
\mathbb{E}\!\left[\ell\big(c_{n+1}; \widehat{\lambda}\big)\right] \;\le\; \alpha,
\label{eq:crc-guarantee}
\end{equation}
so the expected rate at which a hallucination is kept is at most $\alpha$. This is
the harm-blind guarantee we take as our baseline: a single $\lambda$ and a single
$\alpha$ applied to every claim regardless of tier.

\subsection{Severity-aware conformal risk control}
\label{sec:sac}
The marginal guarantee in \eqref{eq:crc-guarantee} averages over tiers, which is
precisely its weakness. Because \textsc{benign} claims are far more numerous than
\textsc{dangerous} ones, the mean loss can sit below $\alpha$ while the loss
restricted to \textsc{dangerous} claims runs well above it: the abundant, easily
filtered benign claims subsidize the dangerous ones. A single budget cannot see
this, because it only ever constrains the pooled rate.

Severity-aware conformal risk control (SAC) removes the pooling. We partition the
calibration set by tier, $\mathcal{C}_t = \{\, i : t(c_i) = t \,\}$ with
$n_t = |\mathcal{C}_t|$, assign each tier its own budget $\alpha_t$, and run the
CRC procedure of \eqref{eq:emp-risk}--\eqref{eq:crc-threshold} \emph{independently
within each tier}:
\begin{equation}
\widehat{\lambda}_t \;=\; \inf\left\{ \lambda :\;
\frac{n_t}{n_t+1}\,\widehat{R}_{n_t}^{\,t}(\lambda) + \frac{B}{n_t+1}
\;\le\; \alpha_t \right\},
\qquad
\widehat{R}_{n_t}^{\,t}(\lambda) = \frac{1}{n_t}\!\sum_{i \in \mathcal{C}_t}
\ell(c_i; \lambda).
\label{eq:sac-threshold}
\end{equation}
At deployment, a claim is kept when its score clears the threshold \emph{of its own
tier}: keep $c$ iff $s(c) \ge \widehat{\lambda}_{t(c)}$. In our experiments we set a
strict dangerous budget and a loose benign budget, $\alpha_{\textsc{dangerous}} =
0.05$ and $\alpha_{\textsc{benign}} = 0.15$, against a marginal baseline of
$\alpha = 0.10$. This is the Mondrian, or group-conditional, construction
\citep{vovk2012conditional} applied to a clinical severity taxonomy: conditioning
the guarantee on the tier rather than reweighting claims inside one budget.

The per-tier guarantee follows immediately, which is the point.

\begin{proposition}[Per-tier risk control]
\label{prop:sac}
Fix a tier $t \in \mathcal{T}$ and suppose the calibration claims
$\{c_i : i \in \mathcal{C}_t\}$ and a fresh test claim of tier $t$ are
exchangeable. Then the tier-$t$ threshold $\widehat{\lambda}_t$ from
\eqref{eq:sac-threshold} satisfies
\begin{equation}
\mathbb{E}\!\left[\ell\big(c; \widehat{\lambda}_t\big) \,\middle|\, t(c) = t\right]
\;\le\; \alpha_t .
\end{equation}
Consequently the bounds hold simultaneously across tiers: the expected rate at
which a \emph{dangerous} hallucination is kept is at most
$\alpha_{\textsc{dangerous}}$, regardless of the tiers' base rates or their
relative sizes.
\end{proposition}

\begin{proof}
Restrict attention to tier $t$. On the subsequence $\mathcal{C}_t$ together with a
test claim of the same tier, the loss $\ell(\cdot; \lambda)$ of \eqref{eq:loss} is
bounded by $B = 1$ and is monotone non-increasing in $\lambda$, and the claims are
exchangeable by assumption. The threshold $\widehat{\lambda}_t$ in
\eqref{eq:sac-threshold} is exactly the conformal-risk-control threshold
\eqref{eq:crc-threshold} computed from the $n_t$ losses of this subsequence.
Applying the conformal risk control theorem \citep{angelopoulos2024crc} to the
subsequence gives $\mathbb{E}[\ell(c; \widehat{\lambda}_t) \mid t(c) = t] \le
\alpha_t$. Since $\mathcal{T}$ partitions the claim space, the tiers are handled by
disjoint calibration subsets and the bounds hold jointly.
\end{proof}

\paragraph{Expected versus per-split risk.}
The guarantee in Proposition~\ref{prop:sac} bounds the \emph{expectation} over the
random draw of calibration and test claims. It does not assert that the realized
tier risk on any single calibration split falls below $\alpha_t$; the empirical
risk fluctuates around its mean from split to split. This distinction drives our
reporting. We measure the realized tier risk as
$\widehat{R}^{\,t}_{\text{test}} = \big(\sum_{c \in t} \mathbf{1}[s(c) \ge
\widehat{\lambda}_t]\,\mathbf{1}[y(c)=1]\big) / |\{c : t(c) = t\}|$ on held-out
claims, and the violation rate as the fraction of calibration/test splits on which
$\widehat{R}^{\,t}_{\text{test}} > \alpha_t$. An expectation-level guarantee is
consistent with a mean that hugs $\alpha_t$ and a violation rate near one half; we
report both rather than the mean alone, so the reader can see where the guarantee
does and does not bite.

\paragraph{Exchangeability and clustering.}
Proposition~\ref{prop:sac} assumes exchangeability at the claim level. Claims are
in fact nested within questions, and claims from the same answer are correlated, so
claim-level exchangeability is an idealization. Rather than assume it away, we
quantify the resulting finite-sample uncertainty with a question-level cluster
bootstrap (Section~\ref{sec:setup-eval}), resampling whole questions so that the
reported confidence intervals reflect the true unit of correlation.

\subsection{Worked example}
\label{sec:example}
The mechanism is easiest to see on a single answer. Suppose a model is asked how to
take metformin and its answer decomposes into four scored, tiered claims:

\begin{center}
\small
\begin{tabular}{llcc}
\toprule
Claim & Tier & $s(c)$ & $y(c)$ \\
\midrule
``Take 5000\,mg per day.'' & \textsc{dangerous} & $0.82$ & $1$ (hallucination) \\
``Avoid metformin in severe kidney disease.'' & \textsc{dangerous} & $0.90$ & $0$ (true) \\
``Metformin derives from French lilac.'' & \textsc{benign} & $0.60$ & $0$ (true) \\
``Metformin was first described in 1922.'' & \textsc{benign} & $0.55$ & $0$ (true) \\
\bottomrule
\end{tabular}
\end{center}

Suppose calibration at the marginal budget $\alpha = 0.10$ returns a single global
threshold $\widehat{\lambda} = 0.70$. Both dangerous claims clear it and are kept,
including the overdose ``5000\,mg'' at $s = 0.82$, while the two true benign claims
are flagged. The harm-blind filter passes the dangerous hallucination because its
score clears a bar set low enough that the abundant, lower-scored benign claims
already hold the pooled rate within budget.

Now apply severity-aware CRC with $\alpha_{\textsc{dangerous}} = 0.05$ and
$\alpha_{\textsc{benign}} = 0.15$. The dangerous tier earns a stricter threshold,
say $\widehat{\lambda}_{\textsc{dangerous}} = 0.88$, and the benign tier a looser
one, $\widehat{\lambda}_{\textsc{benign}} = 0.50$. The overdose claim
($0.82 < 0.88$) is now flagged, the true contraindication ($0.90 \ge 0.88$) is
kept, and both true benign facts ($0.60, 0.55 \ge 0.50$) are surfaced. The scores
and claims are identical; only the allocation of strictness changed. Severity-aware
control catches the dangerous hallucination and retains more true benign
information at the same time, because it stopped spending one budget on two kinds
of error.


\section{Preamble}

The first line of the file must be
\begin{quote}
\begin{verbatim}
\documentclass[11pt]{article}
\end{verbatim}
\end{quote}

To load the style file in the review version:
\begin{quote}
\begin{verbatim}
\usepackage[review]{acl}
\end{verbatim}
\end{quote}
For the final version, omit the \verb|review| option:
\begin{quote}
\begin{verbatim}
\usepackage{acl}
\end{verbatim}
\end{quote}

To use Times Roman, put the following in the preamble:
\begin{quote}
\begin{verbatim}
\usepackage{times}
\end{verbatim}
\end{quote}
(Alternatives like txfonts or newtx are also acceptable.)

Please see the \LaTeX{} source of this document for comments on other packages that may be useful.

Set the title and author using \verb|\title| and \verb|\author|. Within the author list, format multiple authors using \verb|\and| and \verb|\And| and \verb|\AND|; please see the \LaTeX{} source for examples.

By default, the box containing the title and author names is set to the minimum of 5 cm. If you need more space, include the following in the preamble:
\begin{quote}
\begin{verbatim}
\setlength\titlebox{<dim>}
\end{verbatim}
\end{quote}
where \verb|<dim>| is replaced with a length. Do not set this length smaller than 5 cm.

\section{Document Body}

\subsection{Footnotes}

Footnotes are inserted with the \verb|\footnote| command.\footnote{This is a footnote.}

\subsection{Tables and figures}

See Table~\ref{tab:accents} for an example of a table and its caption.
\textbf{Do not override the default caption sizes.}

\begin{table}
  \centering
  \begin{tabular}{lc}
    \hline
    \textbf{Command} & \textbf{Output} \\
    \hline
    \verb|{\"a}|     & {\"a}           \\
    \verb|{\^e}|     & {\^e}           \\
    \verb|{\`i}|     & {\`i}           \\
    \verb|{\.I}|     & {\.I}           \\
    \verb|{\o}|      & {\o}            \\
    \verb|{\'u}|     & {\'u}           \\
    \verb|{\aa}|     & {\aa}           \\\hline
  \end{tabular}
  \begin{tabular}{lc}
    \hline
    \textbf{Command} & \textbf{Output} \\
    \hline
    \verb|{\c c}|    & {\c c}          \\
    \verb|{\u g}|    & {\u g}          \\
    \verb|{\l}|      & {\l}            \\
    \verb|{\~n}|     & {\~n}           \\
    \verb|{\H o}|    & {\H o}          \\
    \verb|{\v r}|    & {\v r}          \\
    \verb|{\ss}|     & {\ss}           \\
    \hline
  \end{tabular}
  \caption{Example commands for accented characters, to be used in, \emph{e.g.}, Bib\TeX{} entries.}
  \label{tab:accents}
\end{table}

As much as possible, fonts in figures should conform
to the document fonts. See Figure~\ref{fig:experiments} for an example of a figure and its caption.

Using the \verb|graphicx| package graphics files can be included within figure
environment at an appropriate point within the text.
The \verb|graphicx| package supports various optional arguments to control the
appearance of the figure.
You must include it explicitly in the \LaTeX{} preamble (after the
\verb|\documentclass| declaration and before \verb|\begin{document}|) using
\verb|\usepackage{graphicx}|.

\begin{figure}[t]
  \includegraphics[width=\columnwidth]{example-image-golden}
  \caption{A figure with a caption that runs for more than one line.
    Example image is usually available through the \texttt{mwe} package
    without even mentioning it in the preamble.}
  \label{fig:experiments}
\end{figure}

\begin{figure*}[t]
  \includegraphics[width=0.48\linewidth]{example-image-a} \hfill
  \includegraphics[width=0.48\linewidth]{example-image-b}
  \caption {A minimal working example to demonstrate how to place
    two images side-by-side.}
\end{figure*}

\subsection{Hyperlinks}

Users of older versions of \LaTeX{} may encounter the following error during compilation:
\begin{quote}
\verb|\pdfendlink| ended up in different nesting level than \verb|\pdfstartlink|.
\end{quote}
This happens when pdf\LaTeX{} is used and a citation splits across a page boundary. The best way to fix this is to upgrade \LaTeX{} to 2018-12-01 or later.

\subsection{Citations}

\begin{table*}
  \centering
  \begin{tabular}{lll}
    \hline
    \textbf{Output}           & \textbf{natbib command} & \textbf{ACL only command} \\
    \hline
    \citep{Gusfield:97}       & \verb|\citep|           &                           \\
    \citealp{Gusfield:97}     & \verb|\citealp|         &                           \\
    \citet{Gusfield:97}       & \verb|\citet|           &                           \\
    \citeyearpar{Gusfield:97} & \verb|\citeyearpar|     &                           \\
    \citeposs{Gusfield:97}    &                         & \verb|\citeposs|          \\
    \hline
  \end{tabular}
  \caption{\label{citation-guide}
    Citation commands supported by the style file.
    The style is based on the natbib package and supports all natbib citation commands.
    It also supports commands defined in previous ACL style files for compatibility.
  }
\end{table*}

Table~\ref{citation-guide} shows the syntax supported by the style files.
We encourage you to use the natbib styles.
You can use the command \verb|\citet| (cite in text) to get ``author (year)'' citations, like this citation to a paper by \citet{Gusfield:97}.
You can use the command \verb|\citep| (cite in parentheses) to get ``(author, year)'' citations \citep{Gusfield:97}.
You can use the command \verb|\citealp| (alternative cite without parentheses) to get ``author, year'' citations, which is useful for using citations within parentheses (e.g. \citealp{Gusfield:97}).

A possessive citation can be made with the command \verb|\citeposs|.
This is not a standard natbib command, so it is generally not compatible
with other style files.

\subsection{References}

\nocite{Ando2005,andrew2007scalable,rasooli-tetrault-2015}

The \LaTeX{} and Bib\TeX{} style files provided roughly follow the American Psychological Association format.
If your own bib file is named \texttt{custom.bib}, then placing the following before any appendices in your \LaTeX{} file will generate the references section for you:
\begin{quote}
\begin{verbatim}
\bibliography{custom}
\end{verbatim}
\end{quote}

You can obtain the complete ACL Anthology as a Bib\TeX{} file from \url{https://aclweb.org/anthology/anthology.bib.gz}.
To include both the Anthology and your own .bib file, use the following instead of the above.
\begin{quote}
\begin{verbatim}
\bibliography{anthology,custom}
\end{verbatim}
\end{quote}

Please see Section~\ref{sec:bibtex} for information on preparing Bib\TeX{} files.

\subsection{Equations}

An example equation is shown below:
\begin{equation}
  \label{eq:example}
  A = \pi r^2
\end{equation}

Labels for equation numbers, sections, subsections, figures and tables
are all defined with the \verb|\label{label}| command and cross references
to them are made with the \verb|\ref{label}| command.

This an example cross-reference to Equation~\ref{eq:example}.

\subsection{Appendices}

Use \verb|\appendix| before any appendix section to switch the section numbering over to letters. See Appendix~\ref{sec:appendix} for an example.

\section{Bib\TeX{} Files}
\label{sec:bibtex}

Unicode cannot be used in Bib\TeX{} entries, and some ways of typing special characters can disrupt Bib\TeX's alphabetization. The recommended way of typing special characters is shown in Table~\ref{tab:accents}.

Please ensure that Bib\TeX{} records contain DOIs or URLs when possible, and for all the ACL materials that you reference.
Use the \verb|doi| field for DOIs and the \verb|url| field for URLs.
If a Bib\TeX{} entry has a URL or DOI field, the paper title in the references section will appear as a hyperlink to the paper, using the hyperref \LaTeX{} package.

\section*{Limitations}

This document does not cover the content requirements for ACL or any
other specific venue.  Check the author instructions for
information on
maximum page lengths, the required ``Limitations'' section,
and so on.

\section*{Acknowledgments}

This document has been adapted
by Steven Bethard, Ryan Cotterell and Rui Yan
from the instructions for earlier ACL and NAACL proceedings, including those for
ACL 2019 by Douwe Kiela and Ivan Vuli\'{c},
NAACL 2019 by Stephanie Lukin and Alla Roskovskaya,
ACL 2018 by Shay Cohen, Kevin Gimpel, and Wei Lu,
NAACL 2018 by Margaret Mitchell and Stephanie Lukin,
Bib\TeX{} suggestions for (NA)ACL 2017/2018 from Jason Eisner,
ACL 2017 by Dan Gildea and Min-Yen Kan,
NAACL 2017 by Margaret Mitchell,
ACL 2012 by Maggie Li and Michael White,
ACL 2010 by Jing-Shin Chang and Philipp Koehn,
ACL 2008 by Johanna D. Moore, Simone Teufel, James Allan, and Sadaoki Furui,
ACL 2005 by Hwee Tou Ng and Kemal Oflazer,
ACL 2002 by Eugene Charniak and Dekang Lin,
and earlier ACL and EACL formats written by several people, including
John Chen, Henry S. Thompson and Donald Walker.
Additional elements were taken from the formatting instructions of the \emph{International Joint Conference on Artificial Intelligence} and the \emph{Conference on Computer Vision and Pattern Recognition}.

% Bibliography entries for the entire Anthology, followed by custom entries
%\bibliography{anthology,custom}
% Custom bibliography entries only
\bibliography{custom}

\appendix

\section{Example Appendix}
\label{sec:appendix}

This is an appendix.

\end{document}
